\documentclass[11pt,a4paper]{article}

\usepackage[margin=1in]{geometry}
\usepackage[T1]{fontenc}
\usepackage{lmodern}
\usepackage{microtype}
\usepackage{booktabs}
\usepackage{array}
\usepackage{amsmath}
\usepackage{amssymb}
\usepackage{hyperref}
\usepackage{enumitem}
\usepackage{xcolor}

\hypersetup{colorlinks=true, linkcolor=blue!50!black,
            urlcolor=blue!60!black, citecolor=blue!50!black}

\newcolumntype{L}[1]{>{\raggedright\arraybackslash}p{#1}}

\title{\vspace{-2em}Response to Sensei's Comments\\
\large How the Paper Was Updated}

\date{}

\begin{document}

\maketitle
\vspace{-3em}

\noindent
This note lists each comment and what was done about it.

% ============================================================================
\section{Question 1: More Fine-Tuning Scenarios}
% ============================================================================

\paragraph{The comment:} The findings came from only one fine-tuning scenario.
Another tuning campaign might behave differently, and showing that would be an
important conclusion.

\paragraph{What I did:} Instead of one retraining, I ran \textbf{15 retrainings}
of C1 (the detector). I changed three things:

\begin{itemize}[noitemsep]
    \item \textbf{How much data:} the full Singapore training set (119 images),
          half of it (60), and a quarter (30).
    \item \textbf{How long we train:} 5, 10, and 20 epochs.
    \item \textbf{Random seed:} every setting repeated 3 times, so we can tell a
          real effect from luck.
\end{itemize}

C2 (prediction) and C3 (planning) were never touched. Only C1 changed. Each run
was measured in the same way as before.

\paragraph{The result:} The entanglement is a range, not one number.

\begin{itemize}[noitemsep]
    \item The retraining \emph{always} moved the driving plan, but by very
          different amounts: from \textbf{1.33~m to 7.80~m}.
    \item The planner got \emph{worse} in 11 runs and \emph{better} in 4.
          So ``retraining hurts the planner'' is not a universal rule.
    \item More fine-tuning data caused a bigger downstream change. The full-data
          runs moved the plan by 6.12~m on average, the smaller sets by about
          3~m.
\end{itemize}

This variety is itself the conclusion: different tuning campaigns really do
produce different entanglements. It is in the paper as \textbf{Finding F6}, with
a full 15-row table.

% ============================================================================
\section{Question 2: Evaluating CARA}
\label{sec:cara}
% ============================================================================

\paragraph{The comment:} CARA was only presented as an idea. Can its
effectiveness be evaluated, at least to some extent?

\paragraph{What CARA is, in one line:} Before we accept a retrained detector,
CARA compares it against the current one and decides: \emph{admit} it, or
\emph{hold} it for inspection.

\paragraph{How I evaluated it:} CARA has a cheap front-end called the
\emph{interface-drift screen}. For each retrained detector, it looks only at that
detector's own outputs and measures how different they are from the old
detector's outputs. It uses four simple things: how many objects are found per
image, how confident the detector is, where the boxes sit, and which classes
appear. These are combined into one \emph{drift score} between 0 and 1.

Table~\ref{tab:drift} shows the full measurement on all 15 runs. It is sorted by
drift score, and the last column marks whether the run really did damage the
planner.

\begin{table}[h!]
\centering
\caption{The complete evaluation data. \emph{Drift} is computed from the
detector's outputs only. \emph{Plan move} is how far the driving plan actually
shifted. \emph{Harmful} means the planner really got worse.}
\label{tab:drift}
\small
\renewcommand{\arraystretch}{1.15}
\begin{tabular}{@{}lrrrl@{}}
\toprule
\textbf{Run} & \textbf{Drift} & \textbf{Plan move (m)} & \textbf{$\Delta_3$} & \textbf{Harmful?} \\
\midrule
full, 10 ep, s1    & 0.582 & 7.80 & $+0.68$ & harmful \\
full, 10 ep, s2    & 0.528 & 7.40 & $+1.34$ & harmful \\
full, 10 ep, s3    & 0.527 & 3.17 & $+0.26$ & harmful \\
full, 20 ep, s2    & 0.525 & 3.17 & $+0.26$ & harmful \\
full, 20 ep, s3    & 0.517 & 3.17 & $+0.26$ & harmful \\
quarter, 10 ep, s2 & 0.478 & 3.29 & $-0.20$ & harmless \\
full, 20 ep, s1    & 0.456 & 2.59 & $+0.24$ & harmful \\
half, 10 ep, s3    & 0.420 & 2.59 & $+0.24$ & harmful \\
half, 10 ep, s2    & 0.409 & 3.17 & $+0.26$ & harmful \\
full, 5 ep, s2     & 0.404 & 1.33 & $-0.00$ & harmless \\
full, 5 ep, s3     & 0.399 & 1.33 & $-0.00$ & harmless \\
quarter, 10 ep, s1 & 0.380 & 3.17 & $+0.26$ & harmful \\
full, 5 ep, s1     & 0.352 & 2.59 & $+0.24$ & harmful \\
quarter, 10 ep, s3 & 0.308 & 3.93 & $-1.04$ & harmless \\
half, 10 ep, s1    & 0.300 & 3.17 & $+0.26$ & harmful \\
\bottomrule
\end{tabular}
\end{table}

I asked two questions of this data.

\paragraph{Evaluation 1: does a high drift score predict a big downstream
problem?} Comparing the Drift column against the Plan-move column gives a rank
correlation of \textbf{0.41}. There is a real signal, but it is loose. It works
best at the top of the table: the two runs that moved the plan the most (7.80~m
and 7.40~m) are also the two highest-drift runs. In the middle of the table the
two columns disagree more often.

\paragraph{Evaluation 2: does it work as an admit/hold decision?} Of the 15 runs,
11 were harmful and 4 were harmless. Reading down the Harmful column shows what a
threshold can and cannot do: most harmful runs sit high, but one harmless run
(quarter, 10 ep, s2) sits at 0.478, and one harmful run (half, 10 ep, s1) sits at
the very bottom at 0.300. Overall the score ranks a harmful run above a harmless
one about \textbf{70\% of the time} (AUROC 0.70). Table~\ref{tab:thresh} shows
the two thresholds.

\begin{table}[h!]
\centering
\caption{Admit/hold decisions at two thresholds.}
\label{tab:thresh}
\small
\renewcommand{\arraystretch}{1.25}
\begin{tabular}{@{}lccccc@{}}
\toprule
\textbf{Threshold} & \textbf{Harmful caught} & \textbf{Harmless admitted} & \textbf{Wrong holds} \\
\midrule
Strict (0.25)   & 11 of 11 (100\%) & 0 of 4 (0\%)  & 4 \\
Balanced (0.41) & 8 of 11 (73\%)   & 3 of 4 (75\%) & 1 \\
\bottomrule
\end{tabular}
\end{table}

\noindent
The strict threshold holds every run. Nothing bad gets through, but nothing good
is ever admitted either, so it is safe and useless at the same time. The balanced
threshold is the usable one: it catches about three-quarters of the harmful runs
while correctly admitting three-quarters of the harmless ones, with only one
wrong hold.

\paragraph{Which signal does the work:} The drift score combines four signals,
and they do not contribute equally. Averaged over the 15 runs, the horizontal box
position matters most (0.81), then the confidence (0.58), then the vertical box
position (0.38), then the number of objects per image (0.36). The class mix
barely matters (0.07). In other words, retraining mainly changes \emph{where} and
\emph{how many} objects are detected, not \emph{which} classes.

\paragraph{Limits of this evaluation:} This is a proof of concept, not a proven
rule. The balanced threshold was chosen using the same 15 runs it is judged on,
so there is no separate test set, and there are only 3 test scenes. The
collision-based part of CARA could not be tested, because none of the runs caused
a collision. All of this is stated in the paper.

Both evaluations appear in the paper, with a figure for Evaluation 1 and a table
for Evaluation 2.

% ============================================================================
\section{The Main New Result: Real Entangled Enhancement}
\label{sec:strict}
% ============================================================================

\paragraph{The problem we had:} Entangled enhancement means: \emph{the upstream
part gets better on its own score, and yet the whole system gets worse.} In all
15 runs above, C1's own score always went \emph{down}. So we could show related
effects, but never the real thing.

\paragraph{Why the score kept dropping:} It was not the pipeline. It was the
data. The 119-image Singapore set is dominated by the \texttt{car} class and the
other classes are rare, so a small fine-tune overfits to cars and the overall
mAP falls.

\paragraph{What I did:} I rebuilt the training set so the classes are more
balanced: images containing rare classes are repeated, so the class mix flattens.
The validation set was left completely untouched, so the scores stay comparable
to every earlier run. Then I ran \textbf{12 more retrainings} (2 balance
strengths $\times$ 2 training lengths $\times$ 3 seeds).

\paragraph{The result:} It worked.

\begin{itemize}[noitemsep]
    \item C1's own score went \textbf{up} in \textbf{6 of the 12} runs. The
          balancing removed the accuracy penalty, as intended.
    \item In \textbf{3 of those runs, the planner got worse at the same time}.
          That is exactly \emph{entangled enhancement}: the detector improved on
          its own metric, and the system still degraded.
\end{itemize}

\begin{table}[h!]
\centering
\caption{The three runs with real entangled enhancement.
$\delta_1$ = detector's own score change (positive = better);
$\Delta_3$ = planner error change (positive = worse).}
\renewcommand{\arraystretch}{1.25}
\begin{tabular}{@{}lrrc@{}}
\toprule
\textbf{Run} & \textbf{$\delta_1$} & \textbf{$\Delta_3$} & \textbf{Entangled enhancement?} \\
\midrule
balance 0.5, 20 epochs, seed 1 & $+0.031$ & $+0.24$ & \textbf{yes} \\
balance 1.0, 20 epochs, seed 3 & $+0.021$ & $+0.24$ & \textbf{yes} \\
balance 1.0, 20 epochs, seed 2 & $+0.012$ & $+0.24$ & \textbf{yes} \\
\bottomrule
\end{tabular}
\end{table}

\noindent
All three happen at the longer training length (20 epochs), under both balance
settings and across two different seeds, so it is not one lucky run.

\paragraph{Why this matters:} Each of these three retrained detectors would have
\emph{passed} a normal component-level check, because the detector's own score
improved. Yet each one made the driving output worse. This is the strongest
argument for why a method like CARA is needed, and it is now the paper's central
empirical claim (\textbf{Finding F7}).

% ============================================================================
\section{Structure and Writing Changes}
% ============================================================================

\paragraph{Paper structure:} The paper now follows the suggested technical-paper
order exactly:

\begin{enumerate}[noitemsep]
    \item Introduction
    \item System Model \emph{(including the definition of entangled enhancement)}
    \item CARA (the proposed technique)
    \item Experiment Setup
    \item Evaluation Results
    \item Discussion
    \item Related Work
    \item Conclusion
\end{enumerate}

CARA was moved out of the back of the paper. It now has its own section
\emph{before} the experiments, so the proposed technique is presented first and
then evaluated. Related Work was moved to the back.

\paragraph{Introduction shortened:} The end-to-end (E2E) discussion had grown too
long. It is now two sentences that simply say why we use a modular pipeline. The
detailed E2E material stayed in Related Work, where it belongs.

\paragraph{Writing:} Two things were fixed throughout the paper:

\begin{itemize}[noitemsep]
    \item \textbf{Long sentences split.} Sentences that joined several different
          topics with ``:'' or ``;'' were divided into shorter ones. For example,
          the sentence about C1, C2 and C3 is now four separate sentences.
    \item \textbf{Terse sentences explained.} Sentences that were hard to follow
          now have an introduction first. For example, ``This paper supplies that
          measurement'' now says \emph{what} measurement before saying it, and
          ``The three isolated-mode rows are exactly zero-delta'' now first
          explains that this is a validity check, not a result.
\end{itemize}


\end{document}
